\documentclass[a4paper, 11pt]{article}

\usepackage[swedish]{babel}
\usepackage[utf8]{inputenc}
\usepackage[T1]{fontenc}
\usepackage{lmodern}

\usepackage{tikz}

\tikzset{
  excl/.style = {circle, draw=black, inner sep=0pt, minimum size=4pt}
}

\tikzset{
  incl/.style = {circle, draw=black, fill=black, inner sep=0pt, minimum size=4pt}
}


\usepackage{pgfplots}


\usepackage{graphicx}
\usepackage{caption}
\usepackage{subcaption}

\usepackage{enumitem}

\usepackage{multicol}

\usepackage{amsmath}


\begin{document}

\title{
    \textbf{MM2001: Inläming 0}
}
\author{Johan Montelius}
\date{HT 2026}

\maketitle



\section*{i)  Persson-Böiers 1.6.1-2}

\subsection*{a)}

Förenkla uttrycket:

$$ \frac{\sqrt[3]{ a^6}}{(\sqrt[3]{\sqrt[4]{a^5}})^6} \times \sqrt{a}$$.


\begin{align*}
  \frac{\sqrt[3]{ a^6}}{(\sqrt[3]{\sqrt[4]{a^5}})^6} \times \sqrt{a} 
  &\;=\; \frac{(a^6)^{1/3}}{(((a^5)^{1/4})^{1/3})^6} \times a^{1/2} & \text{ ty\ } \sqrt[n]{a} = a^{1/n}\\
  &\;=\; \frac{a^{6/3}}{a^{5\times {1/4} \times {1/3} \times 6}} \times a^{1/2} & \text{ ty\ } (a^n)^m = a^{n\times m}\\
  &\;=\; \frac{a^{2} \times a^{1/2} }{a^{5 / 2}}  & \text{utför multiplikationerna och förenklar}\\
  &\;=\; \frac{a^{5/2}}{a^{5/2}}  & \text{ty\ } a^n \times a^m = a^{n +  m} \\
  &\;=\; 1  &  \\
\end{align*}
  

\subsection*{b)}

Visa att $x = 9/4$ är en lösning till ekvationen:

$$x^{x \sqrt{x}} = (x\sqrt{x})^x$$.

\begin{align*}
  (9/4)^{9/4 \sqrt{9/4}}   &\; =\;   (9/4\sqrt{9/4})^{9/4}  & \\
  (9/4)^{9/4 \times 3/2}  &\; =\;   (9/4\times 3/2)^{9/4}  &\; \text{ ty\;} \sqrt{n/m} = \sqrt{n}/\sqrt{m}\\
  (9/4)^{9/4 \times 3/2}  &\; =\;  (9/4)^{9/4} \times  (3/2)^{9/4}   &\; \text{ ty\ } (a \times b)^n = a^n \times b^n\\
  (9/4)^{9/4 \times 1/2}  &\; =\;   (3/2)^{9/4}   &\; \text{ dividera med\ } (9/4)^{9/4}\\
  ((9/4)^{1/2})^{9/4} &\; =\;   (3/2)^{9/4}   &\; \text{ ty\ } a^{n\times m} = (a^n)^m\\
  (3/2)^{9/4} &\; =\;   (3/2)^{9/4}   &\; \text{ ty\ } (a/b)^n = a^n/b^n \\
  \end{align*}

\subsection*{ c)}

Har du nu löst ekvation ovan?

\vspace{10pt}

\noindent Nej - jag har visat att den gäller för $x = 9/4$ men inte löst ut ett
generellet uttryck som ger all värden för $x$. Det kan finnas andra
lösningar. 

Vi kan, genom att utföra mer eller mindre samma operationer, lösa ekvationen. 

\begin{align*}
  (x)^{x \sqrt{x}}   &\; =\;   (x\sqrt{x})^{x}  & \\
  (x)^{x \times \sqrt{x}}  &\; =\;  (x)^{x} \times  (\sqrt{x})^{x}   &\; \text{ ty\ } (a \times b)^n = a^n \times b^n\\
  (\sqrt{x})^{2 \times x \times \sqrt{x}} &\; =\;  (\sqrt{x})^{2\times x} \times (\sqrt{x})^{x}  & \text{skriv om\ } x \text{\ som\ } (\sqrt{x})^2  \\
  (\sqrt{x})^{2 \times x \times \sqrt{x}} &\; =\;  (\sqrt{x})^{2\times x + x} & \text{ ty\ } a^n \times a^m = a^{n+m}\\
  2 \times x \times \sqrt{x} &\; =\;  2\times x + x & \text{ ty\ } a^n = a^m \Rightarrow n = m \\
  2 \times \sqrt{x} &\; =\; 3 & \text{ dividera med\ } x\\
  \sqrt{x} &\; =\; 3/2 & \\
\end{align*}

\noindent Vi får då att vi ar en lösning för all värden på $x$ där $\sqrt{x} =
3/2$. Detta ger naturligtvis att $9/4$ är en lösning men vi har även
$-9/4$ som en lösning.

\section*{ii) Persson-Böiers 1.6.1}

\subsection*{a)}

Yttryck formelrad (27) med ord.

\vspace{10pt}

\noindent  Givet att a är större än noll, att b är större än a och att $\alpha$
är större än noll så gäller att b upphöjt till $\alpha$ är större än a
upphöjt till $\alpha$.


\subsection*{b}

Vad händer i (27) om $\alpha \leq 0$?

\vspace{10pt}

\noindent  Om $\alpha$ tillåts vara noll så är naturligtvis $a^\alpha = b^\alpha$
eftersom vad som helst upphöjt till noll är ett.

  

\end{document}


